\section{研究设计}

\subsection{数据来源与处理}

本文的数据集整合自三个来源。企业财务特征、治理结构和基本经营指标来自国泰安（CSMAR）中国上市公司数据库，该数据库涵盖了沪深A股全部上市公司的年度财务报告信息。供应链相关指标来自万得（Wind）供应链数据库。数字化转型指数的构建基于上市公司年度报告的文本分析，提取数字化相关关键词的词频统计生成。这一方法在近年来中国上市公司数字化转型研究中被广泛采用，其核心逻辑是年报中的数字化词频反映了企业在该领域的战略重视程度和实际投入力度。

初始样本覆盖2015—2024年全部A股上市公司。我们执行了以下筛选程序：剔除金融行业上市公司（其资产负债表结构与实体企业存在本质差异）、剔除ST和*ST企业（其经营状态异常，财务指标不可比）、剔除关键变量缺失的观测值。最终得到包含3,663家上市公司、27,510个企业—年份观测值的非平衡面板数据。为缓解极端值对回归系数的扰动，我们对供应链韧性、流程优化度和盈利能力三个变量在1\%和99\%分位数上进行了缩尾处理。

\subsection{变量定义}

被解释变量。供应链韧性（SCR）综合反映了企业在面临供应链中断冲击时的稳定性和恢复能力，涵盖库存周转效率、供应商集中度规避和供应中断恢复周期等维度。该指数取值为正连续变量，经缩尾后分布于0.233至3.298之间，数值越高表示韧性越强。

核心解释变量。数字化转型指数（DT）基于上市公司年度报告中与数字化相关的关键词词频统计生成，涵盖人工智能、大数据、云计算、区块链、物联网、工业互联网、智能制造、电子商务等数字化技术及其应用场景的相关词汇。该变量的样本均值为38.702，标准差为10.722，表明不同企业间的数字化水平存在显著差异。

中介变量。本文设置三个中介变量以检验传导机制。流程优化度（Process）反映企业在生产运营和供应链管理中流程标准化与效率优化的程度，缩尾后分布于0.378至303.468之间。企业创新能力（Innov）以研发投入强度和专利产出等指标衡量。信息共享水平（InfoShare）为1—4的有序分类变量，反映企业在年报中披露的供应链信息透明度以及与上下游伙伴的信息交互程度。

控制变量。参照既有文献的惯常做法，我们纳入以下控制变量：财务杠杆（Lev），以资产负债率衡量，控制企业资本结构对经营决策的影响；独董比例（Indep），以独立董事占董事会人数的百分比衡量，控制公司治理质量；企业规模（Size），以总资产的自然对数衡量；企业年龄（Age），以上市年限的自然对数衡量；盈利能力（ROA），以总资产收益率衡量。

\subsection{计量模型设定}

\textbf{基准回归模型}。本文采用逐步引入控制变量的双向固定效应模型，以在控制不可观测的企业异质性和时间趋势的同时，观察核心系数在不同控制力度下的稳定性和收敛性：

\begin{equation}
SCR_{it} = \beta_0 + \beta_1 DT_{it} + \gamma X_{it} + \mu_i + \lambda_t + \varepsilon_{it}
\label{eq:baseline}
\end{equation}

其中，$SCR_{it}$为企业$i$在$t$年的供应链韧性，$DT_{it}$为数字化转型指数，$X_{it}$为控制变量向量，$\mu_i$为企业固定效应（吸收不随时间变化的企业特征），$\lambda_t$为年份固定效应（吸收所有企业共同面临的宏观冲击和趋势）。模型从仅包含DT的零控设定（M1）开始，逐步引入财务治理变量（M2）、企业特征变量（M3—M4），至纳入年份固定效应的完整模型（M5/M6）。所有回归在企业层面聚类稳健标准误，以处理序列相关和异方差问题。

\textbf{中介效应检验}。采用\cite{baron1986moderator}的三步法结合手动Sobel检验：

\begin{align}
SCR_{it} &= c \cdot DT_{it} + \gamma X_{it} + \mu_i + \lambda_t + \varepsilon_{it} \label{eq:step1} \\
M_{it} &= a \cdot DT_{it} + \gamma X_{it} + \mu_i + \lambda_t + \varepsilon_{it} \label{eq:step2} \\
SCR_{it} &= c' \cdot DT_{it} + b \cdot M_{it} + \gamma X_{it} + \mu_i + \lambda_t + \varepsilon_{it} \label{eq:step3}
\end{align}

其中$M_{it}$分别取流程优化度、企业创新能力和信息共享水平。间接效应为$a \times b$，其统计显著性通过Sobel检验统计量$z = ab / \sqrt{b^2 SE_a^2 + a^2 SE_b^2}$判定。与常见的Bootstrap检验相比，Sobel检验对正态性假设的依赖较强，但本文的大样本特征（$N=27,510$）使得中心极限定理的适用性得到保障。

\textbf{异质性分析}。采用分组回归和交互项检验两种互补策略。分组回归按企业规模中位数、财务杠杆中位数和疫情发生年份（2020年）将样本一分为二，分别估计并比较子样本间的DT系数差异。交互项检验则在全样本中引入DT与分组虚拟变量的乘积项：$SCR_{it} = \beta_0 + \beta_1 DT_{it} + \beta_2 DT_{it} \times G_{it} + \gamma X_{it} + \mu_i + \lambda_t + \varepsilon_{it}$，$\beta_2$捕捉了分组间的边际效应差异。交互项检验的优势在于能够对系数差异进行正式的统计推断，避免分组回归的比较仅依赖各自系数的显著性而忽略系数间的不确定性。